\mwnewpart
\renewcommand\runningtitle{Section 7. Testing}
\chapter{Testing}
\sectioncontents

\section{Two Emulated PCs}\ixs{testing}

\begin{plainwords}
Every part of RetrOS/LT was run before it was kept, on two kinds of
emulated computer. One is quick and can be driven from a script like a
person at the keyboard; the other is slow but runs IBM's own BIOS of
October 1982 on an emulated 8088, as close to the real Personal Computer as
software gets. Neither is a real Personal Computer, and some things only a
real one can show.
\end{plainwords}

\subsection{QEMU} \file{tests/qpc.py} starts QEMU with diskette and fixed
disk images, types keys through its monitor, and reads the screen from
video memory at B8000h. It can capture the printer port to a file. Its
BIOS is SeaBIOS, not IBM's; the differences that showed are below.

\subsection{MAME's IBM PC 5150} \file{tests/mame5150} runs MAME's model
of the 5150 with IBM's BIOS 1501476 (10/27/82), which the owner supplies;
it is not in the repository. MAME's keyboard model needs a ROM that is not
used, so nothing can be typed: tests run from \file{AUTOEXEC.BAT}, and a
Lua script saves the screen from video memory.

\section{The Tests}

\begin{tabularx}{\textwidth}{@{}P{1.55in}Y@{}}
\toprule
\file{tests/hd\_test.py} & An empty XT disk: \typed{FDISK} (whole disk),
the restart, \typed{FORMAT C:/S/V}, starting from C, and \typed{CHKDSK}.
Ends \shown{HD-TEST OK}.\\
\file{tests/install\_test.py} & An empty XT disk and the two diskettes:
Install from beginning to end, as a person would (the partition, the
restart, the format, the diskette change), then starting from C and
running \typed{VER}, \typed{DEBUG} and \typed{CHKDSK}, reading the
startup files, and \typed{BASIC /ROM} and \typed{BASIC} from the path.
Ends \shown{INSTALL-TEST OK}.\\
\file{tests/basic\_test.py} & The two diskettes: \typed{BASIC /ROM}, then
a BASIC program that checks itself, arithmetic, strings, files,
\typed{PLAY}, the screen, \typed{PAINT}, \typed{GET} and \typed{PUT}, and
error trapping, and writes what failed. Ends \shown{BASIC-TEST OK}.\\
On the 5150 & Starting from 360 and 180 KB diskettes at 256 and 640 KB;
\file{ANSI.SYS}, \typed{MODE}, \typed{GRAPHICS}, \typed{TREE},
\typed{VER} and \typed{CHKDSK} from \file{AUTOEXEC.BAT}; the time from
power on to the date prompt. BASIC at 256 and 640 KB: \typed{BASICA /ROM}
with the stand-in BASIC sockets, and a program of arithmetic, files,
\typed{PLAY}, \typed{PAINT}, \typed{GET} and \typed{PUT} and
\typed{SCREEN}, ending with \typed{SYSTEM}.\\
\bottomrule
\end{tabularx}

Each command was also run through its cases when it was written: every
switch of \typed{FORMAT}, \typed{SYS} onto formatted diskettes,
\typed{DISKCOPY} onto unformatted ones, \typed{DISKCOMP} and \typed{COMP}
on equal and different ones, \typed{ASSIGN} with \typed{FORMAT}'s check,
every form of \typed{MODE}, \typed{GRAPHICS}' printer output turned back
into a picture, \typed{BACKUP} of a 500 KB file and two files of the same
name across three diskettes and \typed{RESTORE} of them, byte for byte,
and each escape sequence of \file{ANSI.SYS}.

\section{The Kernel}

The kernel exerciser of the RetrOS/x86 tests (\file{STRESS.ASM}) was run
on RetrOS/LT on the 5150. Its DOS 2 parts, handles, seeking,
directories, finding files, duplicate handles, file control blocks, memory
allocation, absolute disk reads and the CP/M-style call at offset 5 of the
program segment prefix, all ran. Its later parts call functions of DOS 3,
and interrupt 2Fh, which DOS 2 does not set up; PC DOS 2.10 would fail
there the same way.

\engnote{The CP/M-style call jumps to F01D:FEF0, which on an 8088 wraps
around to 0000:00C0. Under QEMU, whose BIOS leaves the twenty-first
address line on, it does not wrap, and the call runs off into memory above
1 MB. On the 5150 it works. Tests that depend on the wrap are run there.}

\section{The Review}\ixs{review}

When the commands were done, the whole of RetrOS/LT's own source was read
again, line by line, for faults and for things that could be made faster
or smaller. It found:

\begin{longopts}
\lopt{Transfers across 64 KB}{Five programs that call the ROM's disk
  routines directly did not keep their buffers off 64 KB boundaries.
  Shown on \typed{FORMAT}: a resident program sized to make DOS load
  \typed{FORMAT} at segment 0E7Dh put its startup record buffer across one,
  and the old \typed{FORMAT} said \shown{Format failure}; the new one
  formats. (Section~5.)}
\lopt{Starting up}{The startup record read \file{RETBIO.COM} a sector at a
  time, and \file{RETBIO.COM} read \file{RETDOS.COM} a cluster at a time.
  Reading whole tracks and runs of clusters took the time from power on to
  the date prompt, on the 5150 at 640 KB from the 180 KB diskette, from 59
  seconds to 51; most of what remains is the BIOS's memory test.}
\lopt{MODE}{A mistake such as \typed{MODE 80,X} changed the screen mode,
  clearing the screen, before reporting the error.}
\lopt{ANSI.SYS}{\keys{Ctrl}{PrtSc} no longer turned printer echo on and
  off. QEMU's BIOS does not give 7200h for that key, so the test puts 7200h
  in the BIOS's keyboard buffer and reads it through DOS.}
\lopt{FORMAT}{PC DOS 2's \typed{FORMAT} says
  \shown{Formatting...Format complete}; RetrOS/LT's said nothing while it
  worked.}
\lopt{COMP, DISKCOMP}{Their work buffers were part of the file, making
  them 18 KB and 7 KB where IBM's were about 2 KB.}
\end{longopts}

It looked at \typed{ASSIGN}'s changing of the caller's string and kept it
(Section~5). It found nothing wrong in \file{RETBIO.COM}, the
reconstructed console functions, the startup records, \typed{FDISK},
\typed{TREE}, \typed{GRAPHICS}, \typed{BACKUP} or \typed{RESTORE}.

\section{What Has Not Been Tested}

\begin{itemize}
\item \textbf{A real IBM Personal Computer or XT.} Everything here ran in
  emulators. Disk timing, the DMA controller, real diskette drives and
  fixed disk adapters, and keyboards are where an emulator and a real
  machine part.
\item \textbf{A second fixed disk}, and fixed disks other than the XT's
  10 MB geometry.
\item \textbf{Cassette BASIC.} \typed{BASIC /ROM} was seen to refuse
  a computer without it, under QEMU and on the 5150 with empty sockets;
  no IBM BASIC ROM was used, so starting one has not been seen.
\item \textbf{BASIC's serial ports, light pen and joysticks}, which the
  emulators used do not provide in a form the tests could drive.
\item \textbf{Printing under QEMU.} QEMU's BIOS reports a time-out after
  each character it sends to a printer file, so \typed{LPRINT}, like DOS's
  own \code{PRN}, sends the character and reports \shown{Device Timeout}.
  On the 5150 the printer was not tested.
\item \textbf{PC DOS 2.10 itself, side by side.} RetrOS/LT's behavior was
  matched to IBM's documentation and to what is known of PC DOS 2.10; no
  copy of PC DOS was used in making it.
\end{itemize}
